% Causal-Epistemic Type Theory with Mechanized Metatheory
% Target: ICFP 2027
%
% The THEORY paper. Pure formalism + Lean proofs.
% No LOC counts, no benchmarks, no "self-hosted" marketing.
%
% Compile: pdflatex main && bibtex main && pdflatex main && pdflatex main

\documentclass[acmsmall,screen,review,anonymous]{acmart}

% Packages
\usepackage{mathpartir}
\usepackage{stmaryrd}
\usepackage{booktabs}
\usepackage{subcaption}
\usepackage{listings}
\usepackage{xspace}

% Macros
\newcommand{\sounio}{\textsc{Sounio}\xspace}
\newcommand{\know}[1]{\texttt{Knowledge}\langle #1 \rangle}
\newcommand{\knowfull}[4]{\texttt{Knowledge}\langle #1, #2, #3, #4 \rangle}
\newcommand{\cknow}[3]{\texttt{CausalKnowledge}\langle #1, #2, #3 \rangle}
\newcommand{\eps}{\epsilon}
\newcommand{\lean}{\textsc{Lean\,4}\xspace}

% Code listing style
\lstdefinelanguage{sounio}{
  keywords={fn, let, var, if, else, while, for, return, struct, enum, match,
            with, kernel, linear, type, use, effect, handle, perform, do},
  keywordstyle=\color{blue!60!black}\bfseries,
  ndkeywords={Knowledge, CausalKnowledge, CausalGraph, f64, i32, bool,
              IO, Mut, Alloc, GPU, Prob, Epistemic},
  ndkeywordstyle=\color{purple!60!black},
  sensitive=true,
  comment=[l]{//},
  commentstyle=\color{green!40!black}\itshape,
  stringstyle=\color{red!60!black},
  morestring=[b]",
}

\lstset{
  language=sounio,
  basicstyle=\ttfamily\small,
  frame=single,
  numbers=left,
  numberstyle=\tiny\color{gray},
  breaklines=true,
  captionpos=b,
  xleftmargin=2em,
  aboveskip=0.8em,
  belowskip=0.5em,
}

% ACM metadata
\setcopyright{acmcopyright}
\acmJournal{PACMPL}
\acmVolume{1}
\acmNumber{ICFP}
\acmArticle{1}
\acmYear{2027}
\acmMonth{9}
\acmDOI{10.1145/nnnnnnn.nnnnnnn}

\begin{document}

\title{Causal-Epistemic Type Theory: \\
Mechanized Metatheory for Uncertainty-Aware Causal Reasoning}

\author{Anonymous}
\affiliation{%
  \institution{Anonymous Institution}
  \country{Anonymous}
}
\email{anonymous@example.org}

\begin{abstract}
Scientific reasoning requires both \emph{epistemic judgment}---how
precisely is a quantity known?---and \emph{causal judgment}---what
would happen under intervention?
We present a type system that unifies these two concerns.
\emph{Epistemic types} $\know{T}$ attach measurement uncertainty and
provenance to values, with arithmetic that automatically propagates
uncertainty following the Guide to the Expression of Uncertainty in
Measurement (GUM, ISO/IEC~98-3).
\emph{Causal types} $\cknow{Y}{\eps}{\texttt{do}(X)}$ embed Pearl's
do-calculus into the type system, rejecting programs that compute
causal effects from non-identifiable queries.

We prove that the combined system is type-safe (progress and
preservation) and GUM-compliant (the computed uncertainty equals the
GUM combined standard uncertainty for all well-typed programs).
These results compose with linear types (four-modality lattice),
algebraic effects (row-polymorphic effect tracking), and units of
measure (dimensional analysis).

\textbf{All metatheory is mechanized in \lean: 731 definitions and
theorems with zero \texttt{sorry} and zero Mathlib dependency.}
To our knowledge, this is the first mechanized treatment of
Pearl's causal hierarchy in a type system context, and the first
type safety proof for a language combining epistemic uncertainty,
causal reasoning, linear resources, algebraic effects, and
dimensional analysis.
\end{abstract}

\begin{CCSXML}
<ccs2012>
<concept><concept_id>10011007.10011006.10011008.10011009.10011012</concept_id>
<concept_desc>Software and its engineering~Functional languages</concept_desc>
<concept_significance>500</concept_significance></concept>
<concept><concept_id>10003752.10010124.10010138</concept_id>
<concept_desc>Theory of computation~Program verification</concept_desc>
<concept_significance>500</concept_significance></concept>
<concept><concept_id>10003752.10010124.10010131</concept_id>
<concept_desc>Theory of computation~Type theory</concept_desc>
<concept_significance>500</concept_significance></concept>
</ccs2012>
\end{CCSXML}

\ccsdesc[500]{Software and its engineering~Functional languages}
\ccsdesc[500]{Theory of computation~Program verification}
\ccsdesc[500]{Theory of computation~Type theory}

\keywords{epistemic types, causal inference, do-calculus, type safety,
mechanized metatheory, GUM, Lean~4, uncertainty quantification}

\maketitle

% ============================================================================
\section{Introduction}
\label{sec:intro}
% ============================================================================

Every measured quantity carries uncertainty.
A laboratory balance reports mass as $m = 10.023 \pm 0.002\,\text{g}$;
a clinical assay yields serum creatinine as $\text{SCr} = 1.2 \pm 0.1\,\text{mg/dL}$;
an ODE solver accumulates numerical error proportional to step count.
When these quantities enter computation---drug dosing calculations,
climate projections, neural network training---their uncertainties must
propagate correctly through every arithmetic operation,
or the results are scientifically meaningless.

But uncertainty alone is insufficient.
A clinician does not merely ask ``what is this patient's creatinine?''
but ``what \emph{would} the creatinine be if we changed the dosing
regimen?''---a causal question that cannot be answered from
observational data alone without structural assumptions.
Current programming languages offer no type-level support for either
concern: uncertainty is tracked (if at all) by runtime libraries
that provide no static guarantees~\cite{lebigot2024uncertainties,
giordano2016measurements}, and causal reasoning is entirely
outside the type system.

We present a type system that makes both epistemic and causal
reasoning first-class type-level concerns, with the following
contributions:

\begin{enumerate}
  \item A type system for \emph{epistemic values} with formal typing
        rules, operational semantics, and proofs of progress,
        preservation, and GUM compliance soundness
        (\S\ref{sec:epistemic}).

  \item A type system for \emph{causal reasoning} that embeds Pearl's
        do-calculus~\cite{pearl2009causality} into the type system,
        rejecting programs that compute causal effects from
        non-identifiable queries (\S\ref{sec:causal}).

  \item The \emph{composition} of epistemic and causal types with
        linear types, algebraic effects, and units of measure,
        yielding a unified type system for scientific reasoning
        (\S\ref{sec:composition}).

  \item \textbf{Full mechanization in \lean}: 731 definitions and
        theorems covering type safety, epistemic confidence lattices,
        causal type theory with do-calculus, linear types with
        four-modality lattice, algebraic effect rows, and units of
        measure---all with zero \texttt{sorry} and zero Mathlib
        dependency (\S\ref{sec:mechanization}).
\end{enumerate}

% ============================================================================
\section{Epistemic Types}
\label{sec:epistemic}
% ============================================================================

We model epistemic values as first-class typed objects rather than
library wrappers.
The core constructor is $\knowfull{T}{\eps}{\Phi}{\mathit{tv}}$, where
$T$ is the value type, $\eps$ is a standard uncertainty bound, $\Phi$
is provenance, and $\mathit{tv}$ is temporal validity.

\subsection{Syntax and Typing}

\paragraph{Types and terms.}
\[
\begin{array}{rcl}
\tau & ::= & T \mid \knowfull{T}{\eps}{\Phi}{\mathit{tv}}
         \mid \tau_1 \otimes \tau_2 \mid \tau_1 \multimap \tau_2 \mid !\tau \\
T & ::= & \texttt{f64} \mid \texttt{i32} \mid \texttt{bool} \mid \ldots \\
\eps & ::= & r \in \mathbb{R}_{\ge 0} \mid \alpha_\eps \\
\Phi & ::= & \texttt{Measured} \mid \texttt{Computed} \mid \texttt{Derived} \\
\mathit{tv} & ::= & r \in \mathbb{R}_{\ge 0} \mid \infty
\end{array}
\]

The uncertainty propagation target follows the GUM combined standard
uncertainty equation for first-order, uncorrelated inputs:
\begin{equation}
\label{eq:gum}
u_c(y)^2 = \sum_{i=1}^{n} \left(\frac{\partial f}{\partial x_i}\right)^2 u(x_i)^2.
\end{equation}

\paragraph{Core typing rules.}
\[
\inferrule*[right=T-Measure]
  {\Gamma \vdash e : T \\ \eps \in \mathbb{R}_{\ge 0}}
  {\Gamma \vdash \texttt{measure}(e,\eps,\Phi,\mathit{tv}) :
   \knowfull{T}{\eps}{\Phi}{\mathit{tv}}}
\]

\[
\inferrule*[right=T-Add]
  {\Gamma \vdash e_1 : \knowfull{T}{\eps_1}{\Phi_1}{t_1} \\
   \Gamma \vdash e_2 : \knowfull{T}{\eps_2}{\Phi_2}{t_2} \\
   \eps_r = \sqrt{\eps_1^2 + \eps_2^2}}
  {\Gamma \vdash e_1 + e_2 :
   \knowfull{T}{\eps_r}{\texttt{Computed}}{\min(t_1, t_2)}}
\]

\[
\inferrule*[right=T-Mul]
  {\Gamma \vdash e_1 : \knowfull{T}{\eps_1}{\Phi_1}{t_1} \\
   \Gamma \vdash e_2 : \knowfull{T}{\eps_2}{\Phi_2}{t_2} \\
   \eps_r = |v_1 v_2| \sqrt{(\eps_1/v_1)^2 + (\eps_2/v_2)^2}}
  {\Gamma \vdash e_1 \times e_2 :
   \knowfull{T}{\eps_r}{\texttt{Computed}}{\min(t_1, t_2)}}
\]

\[
\inferrule*[right=T-Transform]
  {\Gamma \vdash f : T_1 \times \cdots \times T_n \to T_r \\
   \Gamma \vdash e_i : \knowfull{T_i}{\eps_i}{\Phi_i}{t_i}}
  {\Gamma \vdash \texttt{transform}(f, e_1,\ldots,e_n) :
   \knowfull{T_r}{\eps_r}{\texttt{Computed}}{\min_i t_i}}
\]

\[
\inferrule*[right=T-Extract]
  {\Gamma \vdash e : \knowfull{T}{\eps}{\Phi}{t}}
  {\Gamma \vdash \texttt{extract}(e) : T}
\]

\subsection{Operational Semantics}

Runtime epistemic values are tuples $(v,\eps,\Phi,t)$.
Evaluation is call-by-value with uncertainty propagation embedded in
reduction rules:
\[
\inferrule*[right=E-Add]
  {(v_1,\eps_1,\Phi_1,t_1) + (v_2,\eps_2,\Phi_2,t_2)}
  {(v_1{+}v_2,\sqrt{\eps_1^2+\eps_2^2},\texttt{Computed},\min(t_1,t_2))}
\]
\[
\inferrule*[right=E-Mul]
  {(v_1,\eps_1,\Phi_1,t_1) \times (v_2,\eps_2,\Phi_2,t_2)}
  {(v_1 v_2, |v_1 v_2|\sqrt{(\eps_1/v_1)^2+(\eps_2/v_2)^2}, \texttt{Computed}, \min(t_1,t_2))}
\]
\[
\inferrule*[right=E-Extract]
  {\texttt{extract}(v,\eps,\Phi,t)}
  {v}
\]

Standard $\beta$-reduction and let-binding evaluation are unchanged;
epistemic propagation composes with ordinary CBV semantics.

% ============================================================================
\section{Causal Types}
\label{sec:causal}
% ============================================================================

The causal layer treats structural assumptions as typed program terms.
A \texttt{CausalGraph} value carries nodes, directed edges, and an
acyclicity witness.
This enables static constraints over interventions and adjustment sets.

\subsection{Structural Causal Models}

A causal term is typed as:
\[
\Gamma \vdash G : \texttt{CausalGraph}\langle V,E,\prec \rangle
\]
where $\prec$ is a topological order over $V$.
Edge-level uncertainty is represented epistemically, so graph structure
and confidence are both typed artifacts.

\subsection{Do-Operator at the Type Level}

Interventions are represented by a dedicated type constructor
$\cknow{Y}{\eps}{\texttt{do}(X)}$.
Typing requires identifiability:
\[
\inferrule*[right=T-Do]
  {\Gamma \vdash G : \texttt{CausalGraph} \\
   \texttt{SMT-PROVE}(\texttt{identifiable}(G, X, Y))}
  {\Gamma \vdash \texttt{do}(G, X{=}x) :
   \cknow{Y}{\eps_{\text{causal}}}{\texttt{do}(X)}}
\]
If identifiability cannot be proved, type checking fails.
This makes causal invalidity a compile-time error rather than a
post-hoc warning.

\subsection{d-Separation, Back-Door, and Front-Door Constraints}

Conditional independence constraints are encoded as predicates over
paths in $G$:
\[
\texttt{dsep}(G, A, B \mid Z).
\]
Adjustment-set typing requires either back-door admissibility or
front-door admissibility:
\[
\texttt{SMT-PROVE}(\texttt{backdoor\_valid}(G,X,Y,Z))
\quad \lor \quad
\texttt{SMT-PROVE}(\texttt{frontdoor\_valid}(G,X,Y,Z)).
\]
The implementation in \texttt{stdlib/causal/core.sio} already exposes
\texttt{do\_intervention}, \texttt{backdoor\_adjustment}, and
\texttt{average\_treatment\_effect}; the type theory formalizes these
contracts as static obligations.

\subsection{Causal-Epistemic Interaction}

Interventional outputs remain epistemic:
\[
\Gamma \vdash e : \knowfull{\cknow{Y}{\eps_c}{\texttt{do}(X)}}{\eps_m}{\Phi}{t}.
\]
Here $\eps_c$ captures causal-model uncertainty and $\eps_m$ captures
measurement/estimation uncertainty.
The combined system therefore rules out two failure modes:
(1) unidentifiable intervention queries, and
(2) uncertainty-discarding arithmetic on identified effects.

% ============================================================================
\section{Metatheory}
\label{sec:metatheory}
% ============================================================================

We prove the core soundness properties for the combined calculus.
All statements below are mechanized (\S\ref{sec:mechanization}).

\paragraph{Determinism.}
If $e \to e_1$ and $e \to e_2$, then $e_1 = e_2$.
CBV evaluation plus deterministic uncertainty formulas yields a unique
next state.

\paragraph{Progress.}
If $\emptyset \vdash e : \tau$, then $e$ is a value or there exists
$e'$ such that $e \to e'$.
Intervention terms progress only when identifiability obligations are
discharged.

\paragraph{Preservation.}
If $\Gamma \vdash e : \tau$ and $e \to e'$, then
$\Gamma \vdash e' : \tau$.
Uncertainty annotations remain type-correct across every reduction.

\paragraph{Type safety.}
By progress and preservation, well-typed programs do not get stuck.

\paragraph{GUM soundness.}
For arithmetic terms over uncorrelated epistemic inputs, propagated
uncertainty equals Equation~\eqref{eq:gum}.

\paragraph{Identifiability soundness.}
If $\Gamma \vdash \texttt{do}(G,X{=}x) : \cknow{Y}{\eps}{\texttt{do}(X)}$,
then the identifiability predicate holds in the mechanized causal model.

\paragraph{d-separation monotonicity.}
If $\texttt{dsep}(G,A,B \mid Z)$ and $Z \subseteq Z' \subseteq V$ with
no collider-opening additions, then $\texttt{dsep}(G,A,B \mid Z')$.
This theorem is used to justify conservative adjustment refinement.

% ============================================================================
\section{Composition with Other Type Disciplines}
\label{sec:composition}
% ============================================================================

The causal-epistemic core composes with three additional disciplines.

\paragraph{Linear modalities.}
The four-point modality lattice (Linear, Affine, Relevant,
Unrestricted) constrains resource use.
Epistemic measurements can be assigned Relevant modality to prevent
silent dropping of uncertainty-carrying values.

\paragraph{Algebraic effects.}
Row-polymorphic effects track side effects in typing judgments.
Epistemic operations and intervention operators remain pure at the core
calculus level, while IO/Mut/Prob effects are tracked explicitly in
surface programs.

\paragraph{Units of measure.}
Unit annotations compose orthogonally with epistemic and causal types,
allowing terms such as
$\knowfull{\texttt{f64}[\mathrm{mg/L}]}{\eps}{\Phi}{t}$ and
interventional queries over unit-safe outcomes.
This blocks dimensionally-invalid causal analyses at compile time.

Composition is conservative: adding one discipline does not invalidate
the metatheory of the others, because each extension is proved via
modular lemmas over shared typing/evaluation relations.

% ============================================================================
\section{Mechanization in Lean 4}
\label{sec:mechanization}
% ============================================================================

The full development is mechanized in \lean across 12 modules
(5{,}539 lines, 731 definitions/theorems), with zero
\texttt{sorry} and zero Mathlib dependency.

\begin{table}[t]
\centering
\caption{Lean 4 mechanization inventory.}
\label{tab:lean-inventory}
\begin{tabular}{lrrl}
\toprule
Module & Lines & Defs/Thms & Coverage \\
\midrule
\texttt{SounioSubstitution} & 740 & 83 & Substitution lemma, contexts \\
\texttt{SounioPreservation} & 678 & 80 & Type safety composition \\
\texttt{SounioProgress}     & 557 & 52 & Progress, canonical forms \\
\texttt{SounioEpistemic}    & 505 & 92 & GUM/lattice epistemic proofs \\
\texttt{SounioCausality}    & 562 & 106 & SCMs, do-calculus, d-separation \\
\texttt{SounioEffects}      & 637 & 87 & Effect rows, handlers \\
\texttt{SounioLinear}       & 459 & 64 & Four-modality lattice \\
\texttt{SounioTyping}       & 490 & 42 & Typing judgments \\
\texttt{SounioSemantics}    & 293 & 27 & CBV semantics \\
\texttt{SounioUnits}        & 278 & 64 & Dimensional analysis \\
\texttt{SounioRowPoly}      & 256 & 34 & Row polymorphism \\
\texttt{SounioFormal}       & 84  & 0  & Main import index \\
\midrule
\textbf{Total}              & \textbf{5{,}539} & \textbf{731} & \\
\bottomrule
\end{tabular}
\end{table}

The proof architecture follows Wright--Felleisen
\cite{wright1994type}: preservation and progress are proved
independently and composed into type safety.
The causal module adds mechanized invariants for intervention
idempotence, d-separation properties, and admissible adjustment sets.

Artifact verification entrypoint:
\begin{verbatim}
cd formal/lean4 && lake build
\end{verbatim}

% ============================================================================
\section{Motivating Example}
\label{sec:example}
% ============================================================================

The following sketch is assembled from repository code in
\texttt{examples/real\_world/05\_pkpd\_data\_analysis.sio} and
\texttt{stdlib/causal/core.sio}.
It combines epistemic PK derivation with a typed intervention query:

\begin{lstlisting}[caption={Causal-epistemic pharmacokinetic sketch from repository sources}]
let cl = epistemic_new(27.0, 9.0)      // clearance: 27 +/- 3 L/h
let vd = epistemic_new(77.0, 64.0)     // volume:    77 +/- 8 L
let t_half = epistemic_mul(epistemic_new(0.693, 0.0), epistemic_div(vd, cl))

var dag = dag_new()
dag = dag_add_node(dag, "Ability", NodeType::Confounder)
dag = dag_add_node(dag, "Education", NodeType::Treatment)
dag = dag_add_node(dag, "Income", NodeType::Outcome)
dag = dag_add_edge(dag, "Ability", "Education", beta_new(8.0, 2.0), 0.4, 0.05)
dag = dag_add_edge(dag, "Education", "Income", beta_new(7.0, 3.0), 0.5, 0.08)

let dag_do = do_intervention(dag, "Education", 1.0)
let ate = average_treatment_effect(dag_do, "Education", "Income")
\end{lstlisting}

The key point is semantic, not numerical: half-life carries propagated
measurement uncertainty by construction, while the interventional effect
requires a causally-valid intervention term.
Both obligations are checked in one typed program.

% ============================================================================
\section{Related Work}
\label{sec:related}
% ============================================================================

\paragraph{Uncertainty propagation libraries.}
Python's \texttt{uncertainties}~\cite{lebigot2024uncertainties} and
Julia's \texttt{Measurements.jl}~\cite{giordano2016measurements}
provide practical runtime uncertainty propagation, but they cannot
enforce preservation of uncertainty through static typing.

\paragraph{Probabilistic programming.}
Stan~\cite{carpenter2017stan}, Pyro~\cite{bingham2019pyro}, and
Gen~\cite{cusumano2019gen} model posterior distributions and causal
structures at a richer statistical level.
Our goal is complementary: compile-time guarantees for a first-order
uncertainty and identifiability core inside a general systems language.

\paragraph{Dependent and refinement typing.}
Idris~\cite{brady2013idris} and F*~\cite{swamy2016dependent} can encode
strong program invariants, but do not provide a native calculus for
GUM propagation plus typed interventions.

\paragraph{Units, linearity, and effects.}
Units-of-measure~\cite{kennedy2009units}, linear types
\cite{wadler1990linear}, and algebraic effects
\cite{plotkin2009handlers} each address a major safety axis.
Our contribution is a mechanized account that composes these axes with
epistemic and causal typing in one framework.

\paragraph{Causal inference foundations.}
Pearl's do-calculus~\cite{pearl2009causality} is the semantic basis for
our intervention rules.
We lift identifiability obligations into type checking and mechanize the
resulting meta-theory.

% ============================================================================
\section{Conclusion}
\label{sec:conclusion}
% ============================================================================

We presented a causal-epistemic type theory in which uncertainty-aware
scientific computation and intervention reasoning are both static
language obligations.
Epistemic typing enforces GUM-consistent uncertainty propagation;
causal typing enforces identifiability preconditions for do-operator
queries.
The combined system remains type-safe and composes with linear
modalities, effect rows, and units of measure.

All core claims are backed by a complete \lean mechanization
(731 definitions/theorems, zero \texttt{sorry}).
The resulting foundation supports a pragmatic path toward reproducible
scientific software where numeric confidence and causal validity are
checked before execution, not reconstructed after publication.

\bibliographystyle{ACM-Reference-Format}
\bibliography{references}

\end{document}
